\section{本章小结}
\label{sec:sens-summary}

本章开头问：预算多一点，最优值变化多少，最优解往哪里移？回答这个问题，首先要说清比较的对象。最优值函数$V$是全局的，分支上的值$v$是局部的，解集合$\bm{S}$可能为空或含多个点，只有选定分支或选解规则才有“解的导数”可谈。

最优值的变化有两类结论。约束右端扰动时，零间隙且乘子达到就给出对任意大小的扰动都成立的仿射下界，它不要求凸性，也不要求活动集不变。只有在$V$可微时，下界的斜率才等于实际导数；$V$不可微时，全体最优乘子的相反数构成全部下界斜率，方向导数由其中最不利的一个决定。目标参数的变化由包络关系描述，一般含参约束下导数等于Lagrange函数显含参数的偏导。

最优解的移动由可逆的最优性方程描述。无约束时看Hessian，等式约束时看KKT矩阵，只需在切空间上正定。不等式约束再加上LICQ与严格互补，活动约束便可当作等式处理。隐函数定理给出的是局部结论，定理的条件是充分而非必要的，方程的光滑解也要用二阶条件判定是否极小。

这些局部结论的失效方式彼此不同。活动集改变时，解连续但转折，值函数可能仍然可微。全局最优解切换时，每条分支都光滑，全局值出现折角，解发生跳跃。方程退化时，唯一的解也可能没有有限的导数，而值函数依然可微。判别时要看的是可用的方程、所选的分支和线性方程的可逆性分别是否变化。

计算时，同一次$\bm{K}$分解可以服务多个右端，只要一个输出的梯度时伴随法只需一个右端。隐式导数对应精确的最优性方程，近似解、有限差分和有限次算法展开各自带有不同的误差，预测的有效性要靠余量、乘子符号和小幅重解来核验。第~\ref{chap:simplex-method}章的固定基分析正是这一思路在线性规划上的实例，第~\ref{chap:bilevel-optimization}章将继续用解映射的导数研究双层优化。
